% ===================================================================
%  Dodatek J: Amplituda ogona oscylacyjnego jako stała sprzężenia
%             cząstki z substratem — mechanizm hierarchii mas leptonów
%  TGP v1 — Sesja v33+ (2026-03-27)
% ===================================================================
\section{Amplituda ogona oscylacyjnego i~hierarchia mas leptonów}%
\label{app:ogon-masy}

\begin{quote}
\textbf{Nota kanoniczna (2026-04-07).} Niniejszy dodatek używa Formulacji~B pola TGP: sprzężenie kinetyczne $f(g) = 1 + 4\ln g$, wartość centralna $g_0^e = 1{,}24$, punkt duchowy $g^* = e^{-1/4}$. W~sformułowaniu kanonicznym (Formulacja~A, patrz \texttt{tgp\_companion.tex}) sprzężenie kinetyczne to $K(g) = g^4$, wartość centralna $g_0^e = 0{,}86941$, a~punkt duchowy nie występuje. Obie formulacje dzielą tę samą prawą stronę ODE solitonowego: $g^2(1-g)$. Mechanizm ogona masy ($A_{\rm tail}$) wymaga rekalkulacji w~Formulacji~A; natomiast algebra Koide'a i~równania master F1--F12 są niezależne od wyboru $K(g)$.
\end{quote}

Niniejszy dodatek formalizuje \textbf{Ścieżkę~9} (skrypty \texttt{ex55}--\texttt{ex59}):
mechanizm, w~którym amplituda oscylacyjnego ogona solitonu~$\Phi$ pełni
rolę stałej sprzężenia cząstki z~substratem, a~jej czwarta potęga reprodukuje
stosunek mas mion–elektron $r_{21} = m_\mu/m_e \approx 206{,}77$.

% -------------------------------------------------------------------
\subsection{Profil solitonu i~jego ogon oscylacyjny}\label{app:J-profil}
% -------------------------------------------------------------------

\subsubsection{ODE solitonu TGP}

Profil solitonu radialnie symetrycznego w~TGP spełnia równanie
(por.\ dodatek~\ref{app:substrat}, równanie~\eqref{eq:soliton-ode}):
\begin{equation}\label{eq:J-ode}
  f(g)\,g'' \;+\; \frac{2}{r}\,g' \;=\; V'(g),
\end{equation}
gdzie $g = \Phi/\Phi_0$ jest znormalizowanym polem, $V(g) = g^3/3 - g^4/4$
jest potencjałem TGP z~warunkiem~$\beta = \gamma$ (stwierdzenie~\ref{prop:vacuum-condition}),
oraz:
\begin{equation}\label{eq:J-fg}
  f(g) = 1 + 2\alpha\ln g, \qquad \alpha = 2.
\end{equation}
Punkt osobliwy (,,duch''): $f(g^*) = 0$ przy $g^* = e^{-1/(2\alpha)} \approx 0{,}779$.%
\footnote{\textbf{Nota kanoniczna}: Punkt duchowy $g^*$ jest artefaktem Formulacji~B ($f(g)=1+4\ln g$). W~Formulacji~A (kanonicznej) sprzężenie kinetyczne $K(g) = g^4 > 0$ dla każdego $g > 0$, więc punkt duchowy \emph{nie występuje} i~bariera kinetyczna nie ogranicza spektrum solitonowego.}

Warunki brzegowe: $g(0) = g_0$ (,,strzał''), $g'(0) = 0$ (symetria),
$g(r) \to 1$ dla $r \to \infty$ (próżnia).

\subsubsection{Asymptotyczny ogon oscylacyjny}

Dla $r \to \infty$ soliton zbiega do próżni $g = 1$ z~oscylacyjnym ogonem:
\begin{equation}\label{eq:J-tail}
  g(r) - 1 \;\sim\; \frac{B_{\rm tail}\cos r + C_{\rm tail}\sin r}{r}
  \qquad (r \to \infty),
\end{equation}
gdzie $B_{\rm tail}, C_{\rm tail}$ są stałymi zależnymi od $g_0$.
Amplituda ogona:
\begin{equation}\label{eq:J-Atail}
  A_{\rm tail}(g_0) = \sqrt{B_{\rm tail}^2 + C_{\rm tail}^2},
  \qquad
  \phi_{\rm tail}(g_0) = \mathrm{atan2}(B_{\rm tail},\, C_{\rm tail}).
\end{equation}

\begin{remark}[Fizyczna interpretacja $A_{\rm tail}$]\label{rem:J-Atail-fiz}
W~TGP (Droga~B) każda cząstka fundamentalna tworzy pole przestrzenności
$\sim C_i\,e^{-m_{\rm sp} r}/r$ wokół siebie. W~reżimie ekranowania
$m_{\rm sp} \to 0$ (fala Yukawy przechodzi w~oscylacyjną) ogon~\eqref{eq:J-tail}
reprezentuje \textbf{''rozmycie'' cząstki w~substracie} --- analogon
sprzężenia z~polem kwantowym. Amplituda $A_{\rm tail}$ jest
\textbf{miarą siły tego sprzężenia}: im większa, tym silniej cząstka
oddziałuje z~gęstością wygenerowanej przestrzeni.
\end{remark}

% -------------------------------------------------------------------
\subsection{Warunek kwantowania: czysta faza}\label{app:J-kwant}
% -------------------------------------------------------------------

\begin{proposition}[Warunek kwantowania elektronu]\label{prop:J-quantization}
Spośród ciągłego spektrum $g_0 \in (g^*, \infty)$ wartości wyróżnione
to te, dla których ogon solitonu jest \textbf{czystym sinusem}:
\begin{equation}\label{eq:J-H1}
  B_{\rm tail}(g_0) = 0
  \qquad (\text{hipoteza H1: zerowa faza kosinusowa}).
\end{equation}
Fizycznie: soliton jest w~\textbf{rezonansie z~substratem} ---
amplituda odchodzącej składowej (kosinusowej, fala wychodzącą)
zanika, pozostawiając wyłącznie składową stojącą (sinusową).
Warunek~\eqref{eq:J-H1} jest substratem analogiem warunku
całkowitego odbicia przy impedancji zerowej.
\end{proposition}

\begin{remark}[Weryfikacja numeryczna]\label{rem:J-verif}
Skrypt~\texttt{ex58} (skan gęsty $g_0 \in [1{,}1, 3{,}0]$, $N = 300$):
\begin{itemize}[nosep]
  \item \textbf{Pierwsze zero $B_{\rm tail}$}: $g_0^e = 1{,}2301$
    ($0{,}8\%$ odchylenie od wartości referencyjnej $g_0^e = 1{,}24$).
    Hipoteza~H1 \textbf{potwierdzona dla elektronu}.
  \item \textbf{Drugie zero $B_{\rm tail}$}: $g_0 = 2{,}0797$,
    co daje $(A(2{,}0797)/A(1{,}23))^4 = 375 \neq 206{,}77$.
    Hipoteza~H1 \textbf{nie wyjaśnia} masy mionu przez samo drugie zero.
\end{itemize}
\end{remark}

% -------------------------------------------------------------------
\subsection{Skalowanie $A_\mathrm{tail}$ i~stosunek mas}\label{app:J-skalowanie}
% -------------------------------------------------------------------

\begin{proposition}[Skalowanie $A_{\rm tail} \propto g_0^4$]\label{prop:J-scaling}
W~sektorze bez odbić ($n_{\rm bounce} = 0$) amplituda ogona skala się jak:
\begin{equation}\label{eq:J-scaling}
  A_{\rm tail}(g_0) \;\propto\; g_0^{4{,}12 \pm 0{,}05}
  \qquad (g_0 \in [1{,}1,\, 2{,}2]),
\end{equation}
co jest bliskie $g_0^4$ (wyniki~\texttt{ex57}, fit~log-liniowy,
korekta wykładnika od nieliniowości $f(g)$).
\end{proposition}

\begin{remark}[Kluczowy wynik: stosunek mas z~amplitudy ogona]\label{rem:J-kluczowy}
Przy $g_0^e = 1{,}24$ (elektron) i~$g_0^\mu = 2{,}00$ (mion):
\begin{equation}\label{eq:J-r21-Atail}
  \left(\frac{A_{\rm tail}(g_0^\mu)}{A_{\rm tail}(g_0^e)}\right)^4
  = \left(\frac{A_{\rm tail}(2{,}00)}{A_{\rm tail}(1{,}24)}\right)^4
  = 210{,}1 \;\approx\; 206{,}77 \qquad (\text{odchylenie } 1{,}6\%).
\end{equation}
Wynik pochodzi ze skryptu~\texttt{ex57} (szeroki skan
$g_0 \in [1{,}02, 7{,}0]$, 7/7 testów pozytywnych).
\end{remark}

\begin{hypothesis}[Masa cząstki jako czwarta potęga amplitudy ogona]%
\label{hyp:J-mass-Atail4}
Masa cząstki fundamentalnej w~TGP jest proporcjonalna do czwartej potęgi
amplitudy sprzężenia jej solitonu z~substratem:
\begin{equation}\label{eq:J-mass-hyp}
  m \;\propto\; A_{\rm tail}(g_0)^4.
\end{equation}
\textbf{Uzasadnienie heurystyczne}: w~teorii pola kwantowego amplituda
rozpraszania w~czwartym rzędzie rachunku zaburzeń skaluje się jak
$\mathcal{A}^4$, gdzie $\mathcal{A}$ jest amplitudą sprzężenia.
W~TGP $A_{\rm tail}$ pełni rolę stałej sprzężenia cząstki z~substratem
(uwaga~\ref{rem:J-Atail-fiz}): masa wynikająca z~samoenergetycznej
pętli czwartego rzędu skaluje się jak $A_{\rm tail}^4$.
Status: \statuslabel{Hipoteza} [POST+NUM] (odchylenie $1{,}5\%$);
argument analityczny: $\mathcal{E}^{(2)} = 0$ (tryb zerowy) --- twierdzenie;
jednak $\mathcal{E}^{(3)} \neq 0$ perturbacyjnie (ex148, uw.~\ref{rem:R-E3-status}),
więc przejście $A^2 \to A^4$ wymaga mechanizmu nieperturbacyjnego.
Weryfikacja numeryczna: \texttt{ex142} (12/12 PASS), \texttt{ex143} (10/10 PASS).
Argument heurystyczny: stopień potencjału $V$ (stw.~\ref{prop:K-exponent}).
\end{hypothesis}

% -------------------------------------------------------------------
\begin{proposition}[Masa solitonu TGP: tryb zerowy i~korekta nieliniowa]%
\label{prop:J-zero-mode-mass}
Niech $g(r)$ będzie rozwiązaniem solitonowym ODE TGP
\begin{equation}
  f(g)\,g'' + \frac{2}{r}\,g' = V'(g),
  \qquad f(g) = 1 + 2\alpha\ln g,\quad \alpha = 2,
\end{equation}
z~warunkiem brzegowym $g(r) \to 1$ dla $r \to \infty$.
Niech $\delta(r) = g(r) - 1$ będzie małym odchyleniem w~ogonie.

\medskip\noindent
\textbf{(i) Tryb zerowy zlinearyzowanego ODE.}
Linearyzacja przy $g = 1$ ($f(1) = 1$, $V''(1) = \gamma = 1$) daje:
\begin{equation}
  \delta'' + \frac{2}{r}\,\delta' + \delta = 0.
\end{equation}
Ogólne rozwiązanie regularnie zanikające to $\delta(r) = A_{\rm tail}\,j_0(r) + \ldots$,
gdzie $j_0(r) = \sin(r)/r$ jest sferyczną funkcją Bessela.
Energia kwadratowego członu fluktuacji w~ogonie ($r \to \infty$):
\begin{equation}\label{eq:J-zero-mode-energy}
  \mathcal{E}_{\rm linear}
  = \int_R^\infty \!\!\bigl[(\delta')^2 - \delta^2\bigr]\,4\pi r^2\,dr
  \;\xrightarrow{R \to \infty}\; 0
\end{equation}
(całki kinetyczna i~potencjalna kasują się dokładnie:
$\frac{1}{2}[\sin(2r)]_R^\infty - \frac{1}{2}[\sin(2r)]_R^\infty \equiv 0$).

\medskip\noindent
\textbf{(ii) Konsekwencja dla masy fizycznej.}
Tryb $\delta \propto \sin(r)/r$ jest \emph{trybem zerowym} zlinearyzowanego
funkcjonału energii, tzn.\ jego wkład do całkowitej energii solitonu
pochodzi wyłącznie z~członu nieliniowego (co najmniej czwartego rzędu):
\begin{equation}\label{eq:J-mass-Atail4-analytic}
  \boxed{M_{\rm TGP} = c_M \cdot A_{\rm tail}^4 + \mathcal{O}(A_{\rm tail}^6),}
\end{equation}
gdzie $c_M > 0$ jest stałą wyznaczaną z~całki:
\[
  c_M \;=\; 4\pi \int_0^\infty \mathcal{L}^{(4)}\!\bigl[j_0(r)\bigr]\,r^2\,dr,
\]
a~$\mathcal{L}^{(4)}$ oznacza człon czwartego rzędu w~rozwinięciu
Lagrangianu względem $\delta$.

\medskip\noindent
\textbf{Status}: \statuslabel{Propozycja} [AN($\mathcal{E}^{(2)}$)+POST($A^4$)+NUM]
(kasowanie wirialowe $\mathcal{E}^{(2)}=0$ --- analitycznie;
ale $\mathcal{E}^{(3)} \neq 0$ perturbacyjnie (ex148),
więc krok $A^2 \to A^4$ jest nieperturbacyjny;
tw.~\ref{thm:R-A4}, dod.~\ref{app:zero-mode-A4};
argument heurystyczny: stopień potencjału (stw.~\ref{prop:K-exponent});
weryfikacja: \texttt{ex142} 12/12 PASS, \texttt{ex67} 5/5 PASS).
\end{proposition}
% -------------------------------------------------------------------

% -------------------------------------------------------------------
\subsection{Złota proporcja i~topologia sektorów}\label{app:J-zlota}
% -------------------------------------------------------------------

\begin{remark}[Obserwacja złotej proporcji]\label{rem:J-golden}
Skrypt~\texttt{ex58} ujawnia:
\begin{equation}\label{eq:J-golden}
  \frac{g_0^\mu}{g_0^e} = \frac{2{,}00}{1{,}24} \approx 1{,}613
  \;\approx\; \varphi = \frac{1 + \sqrt{5}}{2} = 1{,}618\ldots
  \qquad (\text{odchylenie } 0{,}3\%).
\end{equation}
Przy dokładnej złotej proporcji: $g_0^\mu = g_0^e \cdot \varphi = 1{,}24 \cdot 1{,}618 = 2{,}006$.
\end{remark}

\begin{hypothesis}[Złota proporcja z~topologii solitonu]\label{hyp:J-golden-topo}
Sektory solitonu $n_{\rm bounce} = 0$ (elektron) i~$n_{\rm bounce} = 1$ (mion)
są topologicznie odrębne --- odpowiadają różnym klasom homotopii trajektorii
$g(r)\colon [0,\infty) \to (g^*, \infty)$ w~przestrzeni konfiguracyjnej
z~barierą kinetyczną $f(g) = 0$ przy $g = g^*$.

Złota proporcja $\varphi$ pojawia się jako wyróżniona relacja między
sąsiednimi sektorami topologicznymi w~układzie z~mnożnikiem Lagrange'a
postaci $\phi^2 = \phi + 1$ (własność samopodobna złotej proporcji:
$g_0^{(n+1)}/g_0^{(n)} \to \varphi$ w~granicy dużych~$n$).

Status: \statuslabel{Hipoteza} topologiczna (wymaga weryfikacji
analitycznej, skrypt~\texttt{ex60}).
\end{hypothesis}

\begin{remark}[Tau jako trzecia generacja]\label{rem:J-tau}
Jeśli $g_0^\mu = g_0^e \cdot \varphi$, to naturalnym kandydatem
na tau jest $g_0^\tau = g_0^\mu \cdot \varphi^k$ dla jakiegoś
$k \geq 1$, co daje:
\[
  r_{31} = \left(\frac{A_{\rm tail}(g_0^\tau)}{A_{\rm tail}(g_0^e)}\right)^4
  \approx (\varphi^k)^{4 \cdot \gamma_{\rm eff}},
\]
gdzie $\gamma_{\rm eff} \approx 4{,}12$ jest wyznaczony z~\texttt{ex57}.
Dla $k = 2$: $r_{31} \approx (g_0^\tau/g_0^e)^{4{,}12}
= (1{,}618^2 \cdot 1{,}24/1{,}24)^{4{,}12} \approx 2{,}618^{4{,}12} \approx 50$
--- dalekie od $r_{31} = 3477$.
Pełne wyprowadzenie hierarchii $r_{31}$ wymaga dodatkowej
zasady selekcji~$g_0^\tau$ (problem otwarty O-J3).
Odwrotna kalkulacja z~warunku $m \propto A_{\rm tail}^4$ (\texttt{ex61})
daje $g_0^\tau \approx 2{,}34$, $m_\tau \approx 1777$~MeV ---
zob.\ hipoteza~\ref{hyp:K-tau} i~tabela~\ref{tab:K-tau-candidates}
w~Dodatku~\ref{app:wkb-atail}.
\end{remark}

% -------------------------------------------------------------------
\subsection{Analityczne szacowanie $A_{\rm tail}$}\label{app:J-analityczne}
% -------------------------------------------------------------------

\subsubsection{Przybliżenie WKB przy barierze kinetycznej}

Dla $g_0$ bliskiego $g^* = e^{-1/4}$ bariera kinetyczna $f(g) \to 0$
tworzy rejon o~dużej "krętności" trajektorii $g(r)$.
W~przybliżeniu WKB (zmienne $s = -\ln g$, $\tilde{r} = r$):
\begin{equation}\label{eq:J-WKB-action}
  A_{\rm tail} \;\sim\; A_0\,\exp\!\left[-\int_{g^*}^{g_0}
    \frac{\sqrt{2|V_{\rm eff}(g)|}}{|f(g)|^{1/2}}\,dg\right],
\end{equation}
gdzie $V_{\rm eff}(g) = V(g) - V(g^*)$ jest potencjałem efektywnym
przesunietym do zera przy barierze.

\begin{remark}[Wykładnik $\approx 4$]\label{rem:J-exp4}
Numeryczna weryfikacja (\texttt{ex57}): $A_{\rm tail} \propto g_0^{4{,}12}$.
Wykładnik $4$ może wynikać z~faktu, że $V(g) = g^3/3 - g^4/4$
jest wielomianem czwartego stopnia z~dominującym ($-g^4/4$) przy $g \to 0$.
Dla małych $g_0$ całka WKB~\eqref{eq:J-WKB-action} skaluje się jak
$\ln(g_0/g^*)^4 + \ldots$, skąd $A_{\rm tail} \propto (g_0/g^*)^4$
(heurystycznie). Argument ze stopnia $V(g)$ i~niezależność wykładnika od~$\alpha$
(skrypt~\texttt{ex62} dla $\alpha \in [1{,}0;\, 3{,}0]$) są sformalizowane
jako twierdzenie~\ref{res:K-SWKB} w~Dodatku~\ref{app:wkb-atail};
pełne wyprowadzenie przez dopasowanie asymptotyczne pozostaje problemem O-J1.
\end{remark}

\subsubsection{Asymptotyczna relacja A-tail z parametrem strzału}

Dla $g_0 \in (1, 2{,}5)$ (sektor $n_{\rm bounce} = 0$) numeryczny fit (\texttt{ex57}):
\begin{equation}\label{eq:J-fit}
  A_{\rm tail}(g_0) \approx c_A \cdot (g_0 - g^*)^{4{,}12},
  \qquad
  c_A \approx 0{,}35 \;\pm\; 0{,}05.
\end{equation}
Wartości odniesienia:
\begin{equation}\label{eq:J-Atail-values}
  A_{\rm tail}(1{,}24) \approx 0{,}287,
  \qquad
  A_{\rm tail}(2{,}00) \approx 1{,}09.
\end{equation}

% -------------------------------------------------------------------
\subsection{Status i~otwarte problemy}\label{app:J-status}
% -------------------------------------------------------------------

\begin{center}\footnotesize
\begin{tabular}{@{}p{4.5cm}lp{4.2cm}p{2.4cm}@{}}
\toprule
\textbf{Zdanie} & \textbf{Status} & \textbf{Warunki / założenia}
  & \textbf{Odsyłacz} \\
\midrule

Ogon oscylacyjny $g(r)-1 \sim (B\cos r + C\sin r)/r$
  & \statuslabel{Twierdzenie}
  & liniowe ODE dla $g-1 \ll 1$ przy $r\to\infty$; $m_{\rm sp}^2 = V''(1) = \gamma$
  & stw.~\ref{prop:N0-6-from-N0-5} \\[3pt]

Skalowanie $A_{\rm tail} \propto g_0^{4{,}12}$
  & \statuslabel{Propozycja}
  & numeryczna weryfikacja \texttt{ex57}; 7/7 PASS
  & stw.~\ref{prop:J-scaling} \\[3pt]

$(A(2{,}00)/A(1{,}24))^4 \approx 206{,}77$ (1{,}6\%)
  & \statuslabel{Propozycja}
  & \texttt{ex57}; $g_0^e, g_0^\mu$ dane
  & uw.~\ref{rem:J-kluczowy} \\[3pt]

$M_{\rm TGP} = c_M A_{\rm tail}^4 + \mathcal{O}(A^6)$ (tryb zerowy)
  & \statuslabel{Propozycja}
  & $\mathcal{E}^{(2)}=0$ [AN]; $\mathcal{E}^{(3)}\neq 0$ [AN/ex148]; krok $A^2 \to A^4$ nieperturbacyjny
  & prop.~\ref{prop:J-zero-mode-mass} \\[3pt]

$m \propto A_{\rm tail}^4$ (hipoteza masy)
  & \statuslabel{Hipoteza}
  & [POST+NUM] ($1{,}6\%$); argument heurystyczny: stopień $V$ (stw.~\ref{prop:K-exponent})
  & hip.~\ref{hyp:J-mass-Atail4} \\[3pt]

$B_{\rm tail}(g_0^e) = 0$ (warunek kwantowania elektronu)
  & \statuslabel{Propozycja}
  & $g_0^e = 1{,}2301$ potwierdzone (\texttt{ex58}; 0{,}8\%)
  & stw.~\ref{prop:J-quantization} \\[3pt]

$g_0^\mu / g_0^e \approx \varphi$ (złota proporcja)
  & \statuslabel{Hipoteza}
  & odchylenie 0{,}3\%; topologiczna motywacja
  & hip.~\ref{hyp:J-golden-topo} \\[3pt]

Analityczne wyprowadzenie $A_{\rm tail}(g_0)$
  & \statuslabel{Propozycja [NUM]}
  & $\nu \approx 3/2$ (WKB turning-point); \texttt{ex218}
  & O-J1 \\[3pt]

Zasada selekcji $g_0^\mu = g_0^e \cdot \varphi$
  & \statuslabel{Hipoteza}
  & $g_0^\mu{=}2{,}00$ vs 2.\ zero $B_{\rm tail}{=}2{,}08$;
    odch.\ $4\%$; sektory w~Dod.~\ref{app:wkb-atail}
  & hip.~\ref{hyp:J-golden-topo}; \texttt{ex60} \\[3pt]

Predykcja $g_0^\tau \approx 2{,}34$, $m_\tau \approx 1777$~MeV
  & \statuslabel{Hipoteza}
  & odwrotna kalk.\ (\texttt{ex61}); Koide $0{,}5\%$
  & hip.~\ref{hyp:K-tau}; dod.~\ref{app:wkb-atail} \\

\bottomrule
\end{tabular}
\end{center}

\paragraph{Otwarte problemy Ścieżki~9.}
\begin{enumerate}[nosep,label=\textbf{O-J\arabic*.}]
  \item \textbf{Analityczne $A_{\rm tail}(g_0)$}: Czy całkę WKB~\eqref{eq:J-WKB-action}
    można zamknąć w~postaci elementarnej? Wykładnik~$4$ odpowiada
    czwartemu stopniowi $V(g)$ (argument sformalizowany w~Dod.~\ref{app:wkb-atail},
    res.~\ref{res:K-SWKB}; potwierdzony numerycznie dla $\alpha\in[1,3]$ przez \texttt{ex62}).
    Pełne zamknięcie analityczne przez metodę dopasowania asymptotycznego:
    program~\ref{prob:K-OJ1}.
  \item \textbf{Zasada selekcji $g_0^\mu$}
    (\statuslabel{Hipoteza}, zob.\ Dod.~\ref{app:wkb-atail}~\S\ref{rem:K-sektory}):
    Warunek $B_{\rm tail}(g_0^\mu) = 0$ daje $g_0^\mu = 2{,}0797 \neq 2{,}00$.
    Co wyróżnia $g_0^\mu = 2{,}00$? Skrypt~\texttt{ex60} potwierdza
    że $g_0^\mu = 2{,}00$ odpowiada granicy sektora topologicznego
    ($n_{\rm cross}$) --- złota proporcja pojawia się jako stosunek
    $z_{\rm next}/z_0$ sąsiednich zer $B_{\rm tail}$, z~odchyleniem~$4\%$
    od dokładnego $\varphi$; analityczne uzasadnienie pozostaje otwarte.
  \item \textbf{Hierarchia $r_{31}$ --- tau}
    (\statuslabel{Hipoteza}, hip.~\ref{hyp:K-tau} w~Dod.~\ref{app:wkb-atail}):
    Odwrotna kalkulacja (\texttt{ex61}) przy $m \propto A_{\rm tail}^4$
    i~$r_{31} = 3477{,}15$ daje $g_0^\tau \approx 2{,}34$,
    $m_\tau \approx 1777$~MeV, zgodne z~Formułą Koidego na poziomie $0{,}5\%$.
    Zasada selekcji (dlaczego akurat $g_0^\tau = 2{,}34$?) pozostaje otwarta.
\end{enumerate}

\begin{remark}[Weryfikacja numeryczna v41: $\varphi$-FP i~bariera ducha]%
\label{rem:J-v41-verification}
Skrypty \texttt{tau\_mass\_selection.py} i~\texttt{tau\_mass\_full\_ode.py}
(v41, 2026-03-30) weryfikują $\varphi$-FP w~obu wariantach ODE.

\medskip\noindent
\textbf{Uproszczone ODE} $f(g)g''+\tfrac{2}{r}g'=V'(g)$
(bez członu $\tfrac{\alpha}{g}(g')^2$):
\begin{itemize}[nosep]
  \item $\varphi$-FP: $g_0^* = 0{,}8993$, $\varphi\cdot g_0^* = 1{,}455$,
    $(A_\mu/A_e)^4 = 206{,}768$ --- \textbf{dokładna zgodność z~PDG}.
  \item Brak bariery ducha --- masa tauonu dostępna.
\end{itemize}

\noindent
\textbf{Pełne ODE} $f(g)[g''+\tfrac{2}{r}g']+\tfrac{\alpha}{g}(g')^2=g^2(1-g)$
(Radau, rtol=$10^{-10}$):
\begin{itemize}[nosep]
  \item $\varphi$-FP: $g_0^* = 0{,}8339$, $\varphi\cdot g_0^* = 1{,}349$,
    $(A_\mu/A_e)^4 = 206{,}768$ --- \textbf{dokładna zgodność z~PDG}.
  \item \textbf{Bariera ducha}: solitony z~$g_0 \gtrsim 1{,}63$ przekraczają
    punkt $g^* = e^{-1/4}$ ($f(g^*)=0$) i~rozwiązanie rozbiera się
    (śledzenie z~50-cyfrową precyzją mpmath potwierdza blowup
    przy $r \approx 3{,}36$ dla $g_0 = 2{,}0$).
  \item Krytyczne wartości: $g_0^{\rm crit} \approx 1{,}63$,
    $g_{\min}(g_0^{\rm crit}) = 0{,}779 \approx g^*$,
    $f(g_{\min}) \approx 0{,}001$.
  \item Mion bezpieczny: $g_{\min}(\varphi\cdot g_0^*) = 0{,}898$,
    margines $g_{\min} - g^* = 0{,}119$.
  \item $\max (A/A_e)^4 \approx 1571$ przy $g_0 = 1{,}613$ ---
    \textbf{stosunek $r_{31} = 3477$ (tauon) nieosiągalny},
    brak czynnika $\approx 2{,}2\times$.
\end{itemize}

\medskip\noindent
\textbf{Wniosek}: Mechanizm $\varphi$-FP jest \textbf{strukturalną własnością}
ODE solitonu TGP --- potwierdzony w~obu wariantach z~identycznym stosunkiem mas.
Pełne ODE obcina jednak spektrum solitonowy barierą ducha: tauon wymaga
\emph{UV-uzupełnienia} ponad $g^*$.

\medskip\noindent
\textbf{ODE substratowe} (UV-uzupełnienie, ghost-free):
\begin{equation}\label{eq:J-substrate-ode}
  g^2\,g'' + g\,(g')^2 + \frac{2}{r}\,g^2\,g' = V'(g),
\end{equation}
z~$K_{\rm sub}(g) = g^2 > 0$ dla każdego $g > 0$ (sek.~\ref{sec:ghost-resolution}).
\begin{itemize}[nosep]
  \item $\varphi$-FP: $g_0^* = 0{,}8694$, $\varphi\cdot g_0^* = 1{,}407$,
    $(A_\mu/A_e)^4 = 206{,}77$ --- zgodność z~PDG ($0{,}04\%$).
  \item \textbf{Tauon osiągalny}: $g_0^\tau = 1{,}729$,
    $r_{31} = (A_\tau/A_e)^4 = 3477{,}2$ --- \textbf{zgodność 0{,}001\% z~PDG}.
  \item $m_\tau^{\rm pred} = 1776{,}83$~MeV (PDG: $1776{,}86$), $\Delta m = -14$~ppm.
  \item Relacja Koidego: $K = 0{,}66666$ (odch.\ $-11$~ppm od dokładnego $2/3$).
  \item Brak bariery ducha --- $g_{\min}(\tau) = 0{,}742$, poniżej $g^*$ bez osobliwości.
\end{itemize}

\medskip\noindent
\begin{tabular}{lccc}
\toprule
  Wielkość & ODE uproszczone & ODE pełne & ODE substratowe \\
\midrule
  $g_0^*$ (elektron) & $0{,}8993$ & $0{,}8339$ & $0{,}8694$ \\
  $\varphi\cdot g_0^*$ (mion) & $1{,}455$ & $1{,}349$ & $1{,}407$ \\
  $(A_\mu/A_e)^4$ & $206{,}768$ & $206{,}768$ & $206{,}77$ \\
  bariera ducha $g_0^{\rm crit}$ & brak & $\approx 1{,}63$ & brak \\
  $g_0^\tau$ & $\approx 2{,}35$ & --- (nieosiągalny) & $1{,}729$ \\
  $r_{31} = (A_\tau/A_e)^4$ & $590$ & --- & $3477{,}2$ \\
  $\Delta r_{31}$ vs PDG & $-83\%$ & --- & $+0{,}001\%$ \\
  $m_\tau$ [MeV] & $302$ & --- & $1776{,}8$ \\
  Koide $K$ & --- & --- & $0{,}66666$ \\
\bottomrule
\end{tabular}

\medskip\noindent
\textbf{Wniosek}: Mechanizm $\varphi$-FP jest strukturalną własnością
\textbf{wszystkich trzech} wariantów ODE.
ODE substratowe~\eqref{eq:J-substrate-ode} stanowi naturalne UV-uzupełnienie
bariery ducha i~reprodukuje \textbf{obie} hierarchie mas ($r_{21}$ i~$r_{31}$)
z~zerową liczbą parametrów wolnych.
Problem O-J3 (zasada selekcji $g_0^\tau$) rozwiązany poniżej
(prop.~\ref{prop:J-koide-chain}).
\end{remark}

\begin{proposition}[Łańcuch $\varphi$-FP + Koide $\to$ trzy masy leptonów]%
\label{prop:J-koide-chain}
Przy założeniach (i)~$m \propto A_{\rm tail}^4$
(hip.~\ref{hyp:J-mass-Atail4}), (ii)~selekcja $\varphi$-FP
($g_0^\mu = \varphi\,g_0^e$), (iii)~formuła Koidego
$K = \frac{m_e+m_\mu+m_\tau}{(\sqrt{m_e}+\sqrt{m_\mu}+\sqrt{m_\tau})^2} = \tfrac{2}{3}$,
\textbf{masy wszystkich trzech naładowanych leptonów wynikają z~ODE substratowego
bez parametrów wolnych}.

\medskip\noindent\textbf{Łańcuch wyprowadzenia:}
\begin{enumerate}[nosep]
\item ODE substratowe~\eqref{eq:J-substrate-ode} $\to$ $A_{\rm tail}(g_0)$ (numerycznie).
\item $\varphi$-FP: wyznacza $g_0^*$ takie, że
      $(A_{\rm tail}(\varphi\,g_0^*)/A_{\rm tail}(g_0^*))^4 = r_{21}$.
      Wynik: $g_0^* = 0{,}8694$, $r_{21} = 206{,}77$.
\item Koide: przy ustalonym $r_{21}$, równanie
      $3(1 + r_{21} + r_{31}) = 2(1 + \sqrt{r_{21}} + \sqrt{r_{31}})^2$
      ma dwa pierwiastki; fizyczny to $r_{31} = 3477{,}4$.
\item Odwrócenie: $g_0^\tau = A_{\rm tail}^{-1}\!\bigl(A_e \cdot r_{31}^{1/4}\bigr) = 1{,}729$.
\end{enumerate}

\medskip\noindent\textbf{Wyniki numeryczne} (\texttt{oj3\_tau\_selection.py}):
\begin{center}\footnotesize
\begin{tabular}{@{}lcccc@{}}
\toprule
  Metoda & $g_0^\tau$ & $r_{31}$ & $m_\tau$ [MeV] & $\Delta$ od PDG \\
\midrule
  PDG (ref.) & --- & $3477{,}15$ & $1776{,}86$ & --- \\
  $\varphi$-FP + Koide & $1{,}7294$ & $3477{,}4$ & $1777{,}0$ & $+60$~ppm \\
  ,,harmoniczna'' $2\,g_0^e$ & $1{,}7389$ & $3715{,}6$ & $1898{,}7$ & $+6{,}9\%$ \\
\bottomrule
\end{tabular}
\end{center}

\medskip\noindent
Hipoteza ,,harmoniczna'' $g_0^\tau = 2\,g_0^e$ daje odchylenie~$6{,}9\%$
i~\textbf{jest wykluczona} jako dokładna zasada.
Formuła Koidego z~$r_{21}(\varphi\text{-FP})$ daje $m_\tau$ z~precyzją
$60$~ppm --- \textbf{kompatybilną z~PDG} na~poziomie niepewności eksperymentalnej
($\sigma_{m_\tau} = 0{,}12$~MeV $\Rightarrow 68$~ppm).

\medskip\noindent
\textbf{Interpretacja}: $\varphi$-FP determinuje \emph{jedną} hierarchię mas ($r_{21}$),
Koide domyka tryplet do dokładnie \emph{trzech} generacji.
Sens fizyczny: Koide jest warunkiem \emph{globalnego domknięcia}
trypletu leptonowego w~sektorze substratowym TGP.

Status: \statuslabel{Propozycja} (warunkowo na hip.~\ref{hyp:J-mass-Atail4}
i~ODE substratowym).
\end{proposition}

\paragraph{Porównanie ze Ścieżkami 1--8.}
Wszystkie poprzednie ścieżki (K*₂/K*₁, M\_conf, argmax, warunek
energetyczny) były zamknięte negatywnie lub osiągały $r_{21}$ tylko
dla wartości $\alpha_K$ niemożliwych przy fizycznym $a_\Gamma$.
Ścieżka~9 jest jakościowo inna: opiera się na~\emph{amplitudzie sprzężenia}
(obserwowalnej: ogon pola), a~nie na~energii solitonu. Odchylenie
$1{,}6\%$ od $206{,}77$ przy $g_0^e = 1{,}24$, $g_0^\mu = 2{,}00$
jest dotychczas \textbf{najlepszym wynikiem} ze wszystkich ścieżek.

\begin{remark}[Precyzyjna weryfikacja numeryczna: 83~ppm]\label{rem:J-precision-ex157}
Niezależna weryfikacja (\texttt{ex157\_one\_param\_prediction.py}, sesja v45)
w~ODE substratowym z~podwyższoną precyzją ($\textsf{rtol} = 10^{-10}$,
$\textsf{atol} = 10^{-12}$) potwierdza łańcuch prop.~\ref{prop:J-koide-chain}:
\begin{center}\footnotesize
\begin{tabular}{@{}lccl@{}}
\toprule
  Wielkość & Predykcja & PDG & $\Delta$ \\
\midrule
  $g_0^e$              & $0{,}869470$ & --- & (bazowy) \\
  $g_0^\mu$            & $1{,}406833$ & --- & $= \varphi\,g_0^e$ \\
  $r_{21}$             & $206{,}77$   & $206{,}77$ & $< 0{,}01\%$ \\[2pt]
  $g_0^\tau$ (K=2/3)   & $1{,}729615$ & --- & \\
  $r_{31}$             & $3477{,}44$  & $3477{,}15$ & $0{,}0083\%$ \\
  $m_\tau$             & $1776{,}97$~MeV & $1776{,}86$~MeV & $0{,}006\%$ \\
\bottomrule
\end{tabular}
\end{center}
Odchylenie $\Delta r_{31} = 0{,}0083\%$ ($\approx 83$~ppm) jest wewnątrz
niepewności eksperymentalnej PDG ($\sigma_{m_\tau} = 0{,}12$~MeV $\Rightarrow 68$~ppm).
Analiza wrażliwości: $\Delta r_{21} = \pm 1\%$ daje $\Delta r_{31} = \pm 0{,}91\%$ (stabilne).

Trzy rozwiązania Koidego $K = 2/3$ istnieją ($g_0^\tau \in \{0{,}785,\; 1{,}186,\; 1{,}729\}$);
wyłącznie trzecie daje fizyczne $r_{31}$ zgodne z~PDG.

\textbf{Niezależne wyprowadzenie masowe (P4, 2026-04-19).} Skrypt
\texttt{research/particle\_sector\_closure/ps2\_g0tau\_first\_principles.py}
wyznacza $g_0^\tau$ \emph{wyłącznie} z~podwójnego warunku masowego
$(m_e, m_\mu) \to (g_0^e, g_0^\mu)$ i~hipotezy $m = c\, A_\text{tail}^4$ (bez
zakładania Koidego). Otrzymane $g_0^{\tau,\text{mass}} = 1{,}73027168$ zgadza
się z~wartością Koide-pochodną $g_0^{\tau,\text{Koide}} = 1{,}72964790$ do
$\sim 0{,}036\%$. To oznacza, że \textbf{Koide $K = 2/3$ jest wyprowadzone,
a~nie założone} --- jedyny fizyczny korzeń $g_0^\tau \approx 1{,}73$ ODE
substratowego jest selekcjonowany przez relację masową $A^4$, a~Koide
wyłania się jako \emph{konsekwencja} (szczegóły:~Dodatek~\ref{app:P4-g0tau}).
\end{remark}

\paragraph{Spójność z~N0 i~duchem teorii.}
Hipoteza~\ref{hyp:J-mass-Atail4} jest zgodna z~założeniami N0 teorii:
\begin{itemize}[nosep]
  \item[N0-1] Masa jest własnością substratową (ogon solitonu),
    nie zewnętrzną stałą Higgsa.
  \item[N0-5] Warunek próżniowy $\beta = \gamma$ determinuje $V(g) = g^3/3 - g^4/4$,
    co daje wykładnik~$\approx 4$ w~skalowaniu $A_{\rm tail}$.
  \item[N0-6] Masa bozonu $m_{\rm sp} = \sqrt{\gamma}$ wyznacza falę oscylacyjną
    ogona (długość fali $\lambda = 2\pi/m_{\rm sp}$).
\end{itemize}

\paragraph{Kontynuacja w~Dodatku~K.}
Analityczne wyniki sesji v34 uzupełniające niniejszy dodatek --- ścisłe
wyprowadzenie formy ogona (prop.~\ref{prop:K-tail-form}),
argument ze stopnia $V(g)$ dla wykładnika~$4$
(res.~\ref{res:K-SWKB}), klasyfikacja sektorów topologicznych
i~predykcja masy tauonu (hip.~\ref{hyp:K-tau}) ---
są zebrane w~Dodatku~\ref{app:wkb-atail}.
